% *****************************************************************************
% **************            SLIDE
% *****************************************************************************
\beamerdefaultoverlayspecification{}
\begin{frame}[fragile]                           % IL FAUT METTRE FRAGILE A CHAQUE FOIS QU'IL Y A DU CODE A METTRE :(
\frametitle{Une courbe}

\begin{center}
\begin{tikzpicture}
\begin{axis}[
    title={Temperature dependence of CuSO\(_4\cdot\)5H\(_2\)O solubility},
    xlabel={Temperature [\textcelsius]},
    ylabel={Solubility [g per 100 g water]},
    xmin=0, xmax=100,
    ymin=0, ymax=120,
    xtick={0,20,40,60,80,100},
    ytick={0,20,40,60,80,100,120},
    legend pos=north west,
    ymajorgrids=true,
    grid style=dashed,
    grid style={redUPEC},
]
\addplot[
    color=blue,
    mark=square,
    ]
    coordinates {
    (0,23.1)(10,27.5)(20,32)(30,37.8)(40,44.6)(60,61.8)(80,83.8)(100,114)
    };
    \legend{CuSO\(_4\cdot\)5H\(_2\)O}
\end{axis}
\end{tikzpicture}
\end{center}

\end{frame}
\beamerdefaultoverlayspecification{<+->}


% *****************************************************************************
% **************            SLIDE
% *****************************************************************************

% Une petite contrainte
% Si vous souhaitez ajouter du texte avec verbatim, il faut absolument faire une commande comme celle-ci
\defverbatim{\textVerbatimOne}{\begin{verbatim}Vous avez déjà de quoi faire avec tout cela !\end{verbatim}}

\beamerdefaultoverlayspecification{}
\begin{frame}[fragile]                         % IL FAUT METTRE FRAGILE A CHAQUE FOIS QU'IL Y A DU VERBATIM :(
\frametitle{Suite}

\visible<1->{
\begin{block}{Un automate}
\begin{tikzpicture}[shorten >=1pt,node distance=1.3cm,on grid,auto, every node/.style={scale=0.75}, scale=0.3] 
   \node[state,initial] (q_0)   {$q_0$}; 
   \node[state] (q_1) [above right=of q_0] {$q_1$}; 
   \node[state] (q_2) [below right=of q_0] {$q_2$}; 
   \node[state,accepting](q_3) [below right=of q_1] {$q_3$};
    \path[->] 
    (q_0) edge  node {1} (q_1)
          edge  node [swap] {1} (q_2)
    (q_1) edge  node  {1} (q_3)
          edge [loop above] node {0} ()
    (q_2) edge  node [swap] {0} (q_3) 
          edge [loop below] node {1} ();
\end{tikzpicture}
\end{block}
}

\visible<2->{
\begin{block}{Un texte avec verbatim}
\textVerbatimOne
\end{block}
}

\end{frame}
\beamerdefaultoverlayspecification{<+->}